% !PRINTED 38
\weberpage{38}
%<<CONT
how, if $S_{\nu-1}$ does not yet exhaust the whole group $S$, one can derive from it a more comprehensive group $S_\nu$ of the same kind.

For every element $\Theta$ of $S$ there will be a certain lowest positive exponent $a_\nu$ for which $\Theta^{a_\nu}$ is contained in $S_{\nu-1}$, and this $a_\nu$ is, by (16), equal to or smaller than $a_{\nu-1}$, and is moreover a power of $p$, since it must be a divisor of the degree of $\Theta$. We choose $\Theta$ so that the exponent $a_\nu$ becomes as large as possible. If $\Theta_1$ is any other element in $S$, and $\Theta_1^h$ the lowest power of $\Theta_1$ that is contained in $S_{\nu-1}$, then $h$ too is a power of $p$ and equal to or smaller than $a_\nu$. Hence $\Theta_1^{a_\nu}$ is contained in $S_{\nu-1}$.

But if we set
\weq{17}{\Theta_1^{a_\nu} = A_1^{\lambda_1} A_2^{\lambda_2} \ldots A_{\nu-1}^{\lambda_{\nu-1}},}
then all the exponents $\lambda$ are divisible by $a_\nu$. For we have
\[
\Theta_1^{a_{\nu-1}} = A_1^{\frac{\lambda_1 a_{\nu-1}}{a_\nu}} A_2^{\frac{\lambda_2 a_{\nu-1}}{a_\nu}} \ldots A_{\nu-1}^{\frac{\lambda_{\nu-1} a_{\nu-1}}{a_\nu}},
\]
and by hypothesis 2) the exponents of $A_1, A_2 \ldots A_{\nu-1}$ on the right-hand side of this formula must be integers divisible by $a_{\nu-1}$. Consequently $\lambda_1, \lambda_2 \ldots \lambda_{\nu-1}$ are divisible by $a_\nu$. We can therefore, returning to the special $\Theta$ considered above and understanding by $h_1, h_2 \ldots h_{\nu-1}$ integers, set
\[
\Theta^{a_\nu} = A_1^{h_1 a_\nu} A_2^{h_2 a_\nu} \ldots A_{\nu-1}^{h_{\nu-1} a_\nu}
\]
and if we then take
\[
A_\nu = \Theta A_1^{-h_1} A_2^{-h_2} \ldots A_{\nu-1}^{-h_{\nu-1}}
\]
then
\weq{18}{A_\nu^{a_\nu} = 1}
is the lowest power of $A_\nu$ that is contained in $S_{\nu-1}$ (since otherwise a still lower power of $\Theta$ would also be contained in $S_{\nu-1}$).

Then
\weq{19}{A_1^{\alpha_1} A_2^{\alpha_2} \ldots A_\nu^{\alpha_\nu}, \quad \alpha_i = 0, 1, 2 \ldots a_i - 1}
is $= 1$ only if $\alpha_1, \alpha_2 \ldots \alpha_\nu$ are divisible by $a_1, a_2 \ldots a_\nu$ and consequently are $= 0$, and accordingly the elements represented in (19) are all different from one another. But these elements form a group $S_\nu$ with the basis $A_1, A_2 \ldots A_\nu$, which satisfies requirement 1).
